\begin{frame}[t]{Calculando uma base de Gröbner}

O algoritmo de Buchberger fornece um procedimento efetivo para transformar um conjunto qualquer de geradores em uma base de Gröbner.

\begin{algbox}{Algoritmo de Buchberger}

\textbf{Entrada:}
\quad $\{g_1,\ldots,g_s\}\subset\mathbb{K}[x_1,\ldots,x_n]$


\medskip

\begin{tabbing}
\qquad\=\qquad\=\qquad\=\kill
\textbf{Defina} $G_0:= \emptyset,G_1:= \{g_1,\ldots,g_s\}$ e inicie $i:=1$\\
\textbf{Enquanto} $G_{i-1} \neq G_i$ \textbf{faça}\\
\> Se existirem $f,h\in G_i$ tais que o resto $r$ da pseudodivisão de $S(f,h)$ por $G_i$ é não nulo\\
\> \> \textbf{Então} $G_{i+1}:=G_i \cup \set{r}$;\\
\> \> \textbf{Senão} $G_{i+1}:=G_i$;\\
\> $i++$;\\
\textbf{retorne} $G$
\end{tabbing}

\textbf{Saída:}
\quad $G:=G_i$ Base de Gröbner para $I$.
\end{algbox}
\end{frame}

\begin{frame}{Demonstração}
    Suponha, por contradição, que ele execute infinitas iterações. Para que o laço não termine, a condição $G_{i-1} \neq G_i$ deve ser satisfeita infinitas vezes, gerando novos restos não nulos.

Considere uma iteração $i$ em que isso ocorre. Seja $r \neq 0$ o resto da pseudodivisão de $S(f,h)$ por $G_i$. Pela propriedade da divisão multivariada, o monômio líder $\mathrm{ml}(r)$ não é divisível por nenhum monômio líder dos elementos de $G_i$. Em particular,\\
\centering \eqbox{$\centering \mathrm{ml}(r)\notin \langle \mathrm{ml}(g)\,;\, g\in G_i\rangle.$}
\pause
\vspace{5mm}
\\\noindent Então, ao definir $G_{i+1}=G_i\cup\{r\}$, obtemos uma inclusão estrita de ideais monomiais:
\centering \eqbox{  $  \centering \langle \mathrm{ml}(G_i)\rangle \subsetneq \langle \mathrm{ml}(G_{i+1})\rangle.$}

\end{frame}

\begin{frame}{Demonstração}
    Como o algoritmo adiciona infinitos restos não nulos, se considerarmos a subsequência de iterações $i_k$ em que o conjunto é efetivamente modificado, produzimos uma cadeia estritamente ascendente infinita de ideais monomiais:\\
\centering \eqbox{  $  \centering \langle \mathrm{ml}(G_{i_1})\rangle \subsetneq \langle \mathrm{ml}(G_{i_2})\rangle \subsetneq \cdots \subsetneq \langle \mathrm{ml}(G_{i_k})\rangle \subsetneq \cdots$}
\vspace{5mm}
\\Um absurdo, já que, pelo Teorema da base de Hilbert, deve existir um $i \in \mathbb{N}$ tal que $\langle \mathrm{ml}(G_i) \rangle = \langle \mathrm{ml}(G_{i+1}) \rangle$. Logo, o algoritmo deve convergir em um número finito de passos. 
\[\qedsymbol\]

\end{frame}